%% 
%% Copyright 2007-2024 Elsevier Ltd
%% 
%% This file is part of the 'Elsarticle Bundle'.
%% ---------------------------------------------
%% 
%% It may be distributed under the conditions of the LaTeX Project Public
%% License, either version 1.3 of this license or (at your option) any
%% later version.  The latest version of this license is in
%%    http://www.latex-project.org/lppl.txt
%% and version 1.3 or later is part of all distributions of LaTeX
%% version 1999/12/01 or later.
%% 
%% The list of all files belonging to the 'Elsarticle Bundle' is
%% given in the file `manifest.txt'.
%% 
%% Template article for Elsevier's document class `elsarticle'
%% with numbered style bibliographic references
%% SP 2008/03/01
%% $Id: elsarticle-template-num.tex 249 2024-04-06 10:51:24Z rishi $
%%
\documentclass[preprint,12pt]{elsarticle}

%% Use the option review to obtain double line spacing
%% \documentclass[authoryear,preprint,review,12pt]{elsarticle}

%% Use the options 1p,twocolumn; 3p; 3p,twocolumn; 5p; or 5p,twocolumn
%% for a journal layout:
%% \documentclass[final,1p,times]{elsarticle}
%% \documentclass[final,1p,times,twocolumn]{elsarticle}
%% \documentclass[final,3p,times]{elsarticle}
%% \documentclass[final,3p,times,twocolumn]{elsarticle}
%% \documentclass[final,5p,times]{elsarticle}
%% \documentclass[final,5p,times,twocolumn]{elsarticle}

%% For including figures, graphicx.sty has been loaded in
%% elsarticle.cls. If you prefer to use the old commands
%% please give \usepackage{epsfig}

%% The amssymb package provides various useful mathematical symbols
\usepackage{amssymb}
%% The amsmath package provides various useful equation environments.
\usepackage{amsmath}
%% The amsthm package provides extended theorem environments
%% \usepackage{amsthm}

%% The lineno packages adds line numbers. Start line numbering with
%% \begin{linenumbers}, end it with \end{linenumbers}. Or switch it on
%% for the whole article with \linenumbers.
%% \usepackage{lineno}

\journal{Journal of Sound and Vibration}

\begin{document}

\begin{frontmatter}

%% Title, authors and addresses

%% use the tnoteref command within \title for footnotes;
%% use the tnotetext command for theassociated footnote;
%% use the fnref command within \author or \affiliation for footnotes;
%% use the fntext command for theassociated footnote;
%% use the corref command within \author for corresponding author footnotes;
%% use the cortext command for theassociated footnote;
%% use the ead command for the email address,
%% and the form \ead[url] for the home page:
%% \title{Title\tnoteref{label1}}
%% \tnotetext[label1]{}
%% \author{Name\corref{cor1}\fnref{label2}}
%% \ead{email address}
%% \ead[url]{home page}
%% \fntext[label2]{}
%% \cortext[cor1]{}
%% \affiliation{organization={},
%%             addressline={},
%%             city={},
%%             postcode={},
%%             state={},
%%             country={}}
%% \fntext[label3]{}

\title{\textbf{Broadband active control of turbulent jet near field by plasma actuators at high Reynolds number} }

%% use optional labels to link authors explicitly to addresses:
%% \author[label1,label2]{}
%% \affiliation[label1]{organization={},
%%             addressline={},
%%             city={},
%%             postcode={},
%%             state={},
%%             country={}}
%%
%% \affiliation[label2]{organization={},
%%             addressline={},
%%             city={},
%%             postcode={},
%%             state={},
%%             country={}}

\author['a']{Victor Kopiev} 
\ead{vkopiev@mktsagi.ru}

\author['a']{Georgy Faranosov}
\author['a']{Vladimir Kopiev}
\author['a']{Oleg Bychkov}%% Author name

\author['b']{Ivan Moralev}%% Author name
\ead{morler@mail.ru}

\author['b']{Alexander Kotvitskii}
\author['b']{Alexander Efimov}

%% Author affiliation
\affiliation['a']{organization={Central Aerohydrodynamic Institute (TsAGI),Moscow Research Complex},%Department and Organization
            addressline={Radio str.,17}, 
            city={Moscow},
            postcode={105005}, 
            country={Russia}}
\affiliation['b']{organization={Joint Institute for High Temperatures (JIHT RAS)},%Department and Organization
            addressline={Izhorskaya str.,13 b.2}, 
            city={Moscow},
            postcode={127415}, 
            country={Russia}}          

            
            


            

%% Abstract
\begin{abstract}
%% Text of abstract
Control of natural instability waves in the free turbulent jet at Re=240-720 k and M=0.25-0.7 was attempted using feed-backward closed-loop system within wave-cancellation approach. Control system was based on barrier discharge plasma actuator, installed on the nozzle lip, real-time processor and microphone. Control law was designed in spectral space using particle swarm algorithm and was based on preliminary linear system identification procedure. Reduction of the pressure pulsations in the jet near field was achieved in the wide frequency band  of Sh=0.3-0.5, where instability waves clearly dominate above non-modal disturbances. Jet excitation experiment reveals details of the plasma actuator operation mechanism.
\end{abstract}

%%Graphical abstract
\begin{graphicalabstract}
%\includegraphics{grabs}
\end{graphicalabstract}

%%Research highlights
\begin{highlights}
\item Research highlight 1
\item Research highlight 2
\end{highlights}

%% Keywords
\begin{keyword}
%% keywords here, in the form: keyword \sep keyword
turbulent jet \sep instability wave \sep plasma actuator \sep closed-loop control 
%% PACS codes here, in the form: \PACS code \sep code

%% MSC codes here, in the form: \MSC code \sep code
%% or \MSC[2008] code \sep code (2000 is the default)

\end{keyword}

\end{frontmatter}

%% Add \usepackage{lineno} before \begin{document} and uncomment 
%% following line to enable line numbers
%% \linenumbers

%% main text
%%

%% Use \section commands to start a section
%% Use \subsubsection, \paragraph, \subparagraph commands to 
%% start 3rd, 4th and 5th level sections.
%% Refer following link for more details.
%% https://en.wikibooks.org/wiki/LaTeX/Document_Structure#Sectioning_commands












\section{Introduction}
\label{sec1} 
%% Labels are used to cross-reference an item using \ref command.

Section text. See Subsection \ref{sec1}.





























\section{Jet model and control strategy}
\label{sec2} 
%% Use \subsection commands to start a subsection.
\subsection{SISO model of the jet}
\label{subsec11}

SISO model of the jet is shown in Figure 1. Jet shear layer is divided into two subsystems by microphones 1 and 2 and the nozzle lip. It is assumed that IW propagation in these sections can be described using transfer functions $P_{1}$ and $P_{2}$.  Note that according to the scheme, $P_{1}$ includes the transfer function of the actuator and its control system. The natural instability waves measured at the position of the first microphone are represented by a disturbance $w$. At the position of the second microphone, we measure the sum of the signal correlated to mic.1, and a non-correlated part $v$. The latter includes the non-linearity on the instability wave developments, non-modal behavior, and instability waves formed due to the distributed shear layer receptivity between the microphones.
\begin{figure}[b]
    \centering
    \includegraphics[width=0.5\linewidth]{fig1.png}
    \caption{SISO model of the turbulent jet}
    \label{fig:1}
\end{figure}

%% https://en.wikibooks.org/wiki/LaTeX/Importing_Graphics#Importing_external_graphics


System identification is performed in open loop mode for the forced and unforced cases. The values, estimated in the unforced case, are marked below by the upper index $^0$ and for the forced by the upper index $^1$.
In the natural case, we obtain the cross-spectral matrix for microphones 1,2:
\begin{equation} \label{eq:1}
    S^0 =
    \begin{pmatrix}
        S_{11}^0 &S_{12}^0 \\
        S_{21}^0 & S_{22}^0
    \end{pmatrix}=
    \begin{pmatrix}
            S_{ww} & P_{2}S_{ww} \\ 
           c.c. & P_{2}P_{2}^*S_{11}^0+S_{vv} 
    \end{pmatrix}
\end{equation}
System identification was fullfilled by forcing the jet by a wide-band noise, nearly flat in the range 100-2000 Hz. Noise samples were generated numerically and translated from the PC using a sound card, with the control signal registered together with the signals from microphones. The noise signal was used to modulate the HF voltage amplitude in the actuator. In the forced case, the relations for both the microphones and the input signal run as follows:
\begin{eqnarray} \label{eq:2}
    &S^1 =
    \begin{pmatrix}
        S_{11}^1 &S_{12}^1 &S_{1u}^1 \\
        S_{u1}^1 &S_{u2}^1 &S_{uu}^1 \\
        \end{pmatrix}= \nonumber    \\
   \nonumber \\
    &\begin{pmatrix}
        S_{ww}+P_{1}P_{1}^*S_{uu} & S_{ww}+P_{2}P_{1}P_{1}^*S_{uu} & P_{1}S_{uu}\\
        c.c. & P_{2}S_{ww} +P_{1}P_{1}^*S_{uu} & P_{2}P_{2}^*S_{11}^0+S_{vv}  \\
        c.c. & c.c. & S_{uu}
    \end{pmatrix}
\end{eqnarray}
Resolving these two equations, we determine the parameters of our system
\begin{equation} \label{eq:3}
    \begin{matrix}
        S_{ww}=S_{11}^0      &&    S_{vv}=S_{22}^0-P_{2}P_{2}^*S_{11}^0 \\
        P_{2}=S_{12}^0/S_{11}^0  &&    P_{1}=S_{1u}^1/S_{uu}^1
    \end{matrix}
\end{equation}

\section{Experimental setup} \label{sec3}
\subsubsection{Mathematics}
%% Inline mathematics is tagged between $ symbols.
This is an example for the symbol $\alpha$ tagged as inline mathematics.

%% Displayed equations can be tagged using various environments. 
%% Single line equations can be tagged using the equation environment.
\begin{equation}
f(x) = (x+a)(x+b)
\end{equation}

%% Unnumbered equations are tagged using starred versions of the environment.
%% amsmath package needs to be loaded for the starred version of equation environment.
\begin{equation*}
f(x) = (x+a)(x+b)
\end{equation*}

%% align or eqnarray environments can be used for multi line equations.
%% & is used to mark alignment points in equations.
%% \\ is used to end a row in a multiline equation.
\begin{align}\label{eq:1}
 f(x) &= (x+a)(x+b) \\
      &= x^2 + (a+b)x + ab
\end{align}

\begin{eqnarray}
 f(x) &=& (x+a)(x+b) \nonumber\\ %% If equation numbering is not needed for a row use \nonumber.
      &=& x^2 + (a+b)x + ab
\end{eqnarray}

%% Unnumbered versions of align and eqnarray
\begin{align*}
 f(x) &= (x+a)(x+b) \\
      &= x^2 + (a+b)x + ab
\end{align*}

\begin{eqnarray*}
 f(x)&=& (x+a)(x+b) \\
     &=& x^2 + (a+b)x + ab
\end{eqnarray*}
\eqref{eq:1}
%% Refer following link for more details.
%% https://en.wikibooks.org/wiki/LaTeX/Mathematics
%% https://en.wikibooks.org/wiki/LaTeX/Advanced_Mathematics

%% Use a table environment to create tables.
%% Refer following link for more details.
%% https://en.wikibooks.org/wiki/LaTeX/Tables
\begin{table}[t]%% placement specifier
%% Use tabular environment to tag the tabular data.
%% https://en.wikibooks.org/wiki/LaTeX/Tables#The_tabular_environment
\centering%% For centre alignment of tabular.
\begin{tabular}{l c r}%% Table column specifiers
%% Tabular cells are separated by &
  1 & 2 & 3 \\ %% A tabular row ends with \\
  4 & 5 & 6 \\
  7 & 8 & 9 \\
\end{tabular}
%% Use \caption command for table caption and label.
\caption{Table Caption}\label{fig1}
\end{table}


%% Use figure environment to create figures
%% Refer following link for more details.
%% https://en.wikibooks.org/wiki/LaTeX/Floats,_Figures_and_Captions
\begin{figure}[t]%% placement specifier
%% Use \includegraphics command to insert graphic files. Place graphics files in 
%% working directory.
\centering%% For centre alignment of image.
\includegraphics{example-image-b}
%% Use \caption command for figure caption and label.
\caption{Figure Caption}\label{fig1}
%% https://en.wikibooks.org/wiki/LaTeX/Importing_Graphics#Importing_external_graphics
\end{figure}


%% The Appendices part is started with the command \appendix;
%% appendix sections are then done as normal sections
\appendix
\section{Example Appendix Section}
\label{app1}

Appendix text.

%% For citations use: 
%%       \cite{<label>} ==> [1]

%%
Example citation, See \cite{lamport94}.

%% If you have bib database file and want bibtex to generate the
%% bibitems, please use
%%
%%  \bibliographystyle{elsarticle-num} 
%%  \bibliography{<your bibdatabase>}

%% else use the following coding to input the bibitems directly in the
%% TeX file.

%% Refer following link for more details about bibliography and citations.
%% https://en.wikibooks.org/wiki/LaTeX/Bibliography_Management

\begin{thebibliography}{00}

%% For numbered reference style
%% \bibitem{label}
%% Text of bibliographic item

\bibitem{lamport94}
  Leslie Lamport,
  \textit{\LaTeX: a document preparation system},
  Addison Wesley, Massachusetts,
  2nd edition,
  1994.

\end{thebibliography}
\end{document}

\endinput
%%
%% End of file `elsarticle-template-num.tex'.
